%%% This file is automatically generated by LeanArchitect. %%%

%%% Macro definitions for \inputleannode, \inputleanmodule %%%

\makeatletter

% \newleannode{name}{latex} defines a new Lean node
\providecommand{\newleannode}[2]{%
  \expandafter\gdef\csname leannode@#1\endcsname{#2}}
% \inputleannode{name} inputs a Lean node
\providecommand{\inputleannode}[1]{%
  \csname leannode@#1\endcsname}

% \newleanmodule{module}{latex} defines a new Lean module
\providecommand{\newleanmodule}[2]{%
  \expandafter\gdef\csname leanmodule@#1\endcsname{#2}}
% \inputleanmodule{module} inputs a Lean module
\providecommand{\inputleanmodule}[1]{%
  \csname leanmodule@#1\endcsname}

\makeatother

%%% Start of main content %%%

\newleannode{zeroTendstoDiff}{\input{/home/runner/work/trough/trough/library/.lake/build/blueprint/module/PrimeNumberTheoremAnd/PerronFormula.artifacts/zeroTendstoDiff}}

\newleannode{RectangleIntegral_tendsTo_VerticalIntegral}{\input{/home/runner/work/trough/trough/library/.lake/build/blueprint/module/PrimeNumberTheoremAnd/PerronFormula.artifacts/RectangleIntegral_tendsTo_VerticalIntegral}}

\newleannode{RectangleIntegral_tendsTo_UpperU}{\input{/home/runner/work/trough/trough/library/.lake/build/blueprint/module/PrimeNumberTheoremAnd/PerronFormula.artifacts/RectangleIntegral_tendsTo_UpperU}}

\newleannode{RectangleIntegral_tendsTo_LowerU}{\input{/home/runner/work/trough/trough/library/.lake/build/blueprint/module/PrimeNumberTheoremAnd/PerronFormula.artifacts/RectangleIntegral_tendsTo_LowerU}}

\newleannode{limitOfConstant}{\input{/home/runner/work/trough/trough/library/.lake/build/blueprint/module/PrimeNumberTheoremAnd/PerronFormula.artifacts/limitOfConstant}}

\newleannode{limitOfConstantLeft}{\input{/home/runner/work/trough/trough/library/.lake/build/blueprint/module/PrimeNumberTheoremAnd/PerronFormula.artifacts/limitOfConstantLeft}}

\newleannode{tendsto_rpow_atTop_nhds_zero_of_norm_lt_one}{\input{/home/runner/work/trough/trough/library/.lake/build/blueprint/module/PrimeNumberTheoremAnd/PerronFormula.artifacts/tendsto_rpow_atTop_nhds_zero_of_norm_lt_one}}

\newleannode{tendsto_rpow_atTop_nhds_zero_of_norm_gt_one}{\input{/home/runner/work/trough/trough/library/.lake/build/blueprint/module/PrimeNumberTheoremAnd/PerronFormula.artifacts/tendsto_rpow_atTop_nhds_zero_of_norm_gt_one}}

\newleannode{isHolomorphicOn}{\input{/home/runner/work/trough/trough/library/.lake/build/blueprint/module/PrimeNumberTheoremAnd/PerronFormula.artifacts/isHolomorphicOn}}

\newleannode{integralPosAux}{\input{/home/runner/work/trough/trough/library/.lake/build/blueprint/module/PrimeNumberTheoremAnd/PerronFormula.artifacts/integralPosAux}}

\newleannode{vertIntBound}{\input{/home/runner/work/trough/trough/library/.lake/build/blueprint/module/PrimeNumberTheoremAnd/PerronFormula.artifacts/vertIntBound}}

\newleannode{vertIntBoundLeft}{\input{/home/runner/work/trough/trough/library/.lake/build/blueprint/module/PrimeNumberTheoremAnd/PerronFormula.artifacts/vertIntBoundLeft}}

\newleannode{isIntegrable}{\input{/home/runner/work/trough/trough/library/.lake/build/blueprint/module/PrimeNumberTheoremAnd/PerronFormula.artifacts/isIntegrable}}

\newleannode{tendsto_zero_Lower}{\input{/home/runner/work/trough/trough/library/.lake/build/blueprint/module/PrimeNumberTheoremAnd/PerronFormula.artifacts/tendsto_zero_Lower}}

\newleannode{Perron.tendsto_zero_Upper}{\input{/home/runner/work/trough/trough/library/.lake/build/blueprint/module/PrimeNumberTheoremAnd/PerronFormula.artifacts/Perron.tendsto_zero_Upper}}

\newleannode{formulaLtOne}{\input{/home/runner/work/trough/trough/library/.lake/build/blueprint/module/PrimeNumberTheoremAnd/PerronFormula.artifacts/formulaLtOne}}

\newleannode{keyIdentity}{\input{/home/runner/work/trough/trough/library/.lake/build/blueprint/module/PrimeNumberTheoremAnd/PerronFormula.artifacts/keyIdentity}}

\newleannode{diffBddAtZero}{\input{/home/runner/work/trough/trough/library/.lake/build/blueprint/module/PrimeNumberTheoremAnd/PerronFormula.artifacts/diffBddAtZero}}

\newleannode{diffBddAtNegOne}{\input{/home/runner/work/trough/trough/library/.lake/build/blueprint/module/PrimeNumberTheoremAnd/PerronFormula.artifacts/diffBddAtNegOne}}

\newleannode{residueAtZero}{\input{/home/runner/work/trough/trough/library/.lake/build/blueprint/module/PrimeNumberTheoremAnd/PerronFormula.artifacts/residueAtZero}}

\newleannode{residueAtNegOne}{\input{/home/runner/work/trough/trough/library/.lake/build/blueprint/module/PrimeNumberTheoremAnd/PerronFormula.artifacts/residueAtNegOne}}

\newleannode{residuePull1}{\input{/home/runner/work/trough/trough/library/.lake/build/blueprint/module/PrimeNumberTheoremAnd/PerronFormula.artifacts/residuePull1}}

\newleannode{residuePull2}{\input{/home/runner/work/trough/trough/library/.lake/build/blueprint/module/PrimeNumberTheoremAnd/PerronFormula.artifacts/residuePull2}}

\newleannode{contourPull3}{\input{/home/runner/work/trough/trough/library/.lake/build/blueprint/module/PrimeNumberTheoremAnd/PerronFormula.artifacts/contourPull3}}

\newleannode{formulaGtOne}{\input{/home/runner/work/trough/trough/library/.lake/build/blueprint/module/PrimeNumberTheoremAnd/PerronFormula.artifacts/formulaGtOne}}

\newleanmodule{PrimeNumberTheoremAnd.PerronFormula}{
In this section, we prove the Perron formula, which plays a key role in our proof of Mellin
inversion.


The following is preparatory material used in the proof of the Perron formula, see Lemma
\ref{formulaLtOne}.


\inputleannode{zeroTendstoDiff}

\inputleannode{RectangleIntegral_tendsTo_VerticalIntegral}

\inputleannode{RectangleIntegral_tendsTo_UpperU}

\inputleannode{RectangleIntegral_tendsTo_LowerU}

TODO : Move to general section


\inputleannode{limitOfConstant}

\inputleannode{limitOfConstantLeft}

\inputleannode{tendsto_rpow_atTop_nhds_zero_of_norm_lt_one}

\inputleannode{tendsto_rpow_atTop_nhds_zero_of_norm_gt_one}

\inputleannode{isHolomorphicOn}

\inputleannode{integralPosAux}

\inputleannode{vertIntBound}

\inputleannode{vertIntBoundLeft}

\inputleannode{isIntegrable}

\inputleannode{tendsto_zero_Lower}

\inputleannode{Perron.tendsto_zero_Upper}

We are ready for the first case of the Perron formula, namely when $x<1$:


\inputleannode{formulaLtOne}

The second case is when $x>1$.
Here are some auxiliary lemmata for the second case.
TODO: Move to more general section


\inputleannode{keyIdentity}

\inputleannode{diffBddAtZero}

\inputleannode{diffBddAtNegOne}

\inputleannode{residueAtZero}

\inputleannode{residueAtNegOne}

\inputleannode{residuePull1}

\inputleannode{residuePull2}

\inputleannode{contourPull3}

\inputleannode{formulaGtOne}

The two together give the Perron formula. (Which doesn't need to be a separate lemma.)

For $x>0$ and $\sigma>0$, we have
$$
\frac1{2\pi i}
\int_{(\sigma)}\frac{x^s}{s(s+1)}ds = \begin{cases}
1-\frac1x & \text{ if }x>1\\
0 & \text{ if } x<1
\end{cases}.
$$

}